\section{The centred limit and the proof of \texorpdfstring{\cref{thm:main}}{Theorem 1.1}}\label{sec:limit}

We keep $0<q<1$. The terms $h_N(k)$ are largest near the middle $k\approx N$ of
their range, so we re-index around it. For every integer $r$ put
\begin{equation}\label{eq:br}
b_r=\frac{q^{2r(r+1)}}{\poch{q}{Q}{r+1}^2},
\end{equation}
using the negative-index convention \eqref{eq:negindex} when $r+1<0$. Unwinding
that convention gives, for $j\ge0$,
\begin{equation}\label{eq:bneg}
b_{-j-1}=q^{2j}\,\poch{q}{Q}{j}^2 .
\end{equation}
For integers $N\ge1$ and $-N\le r\le N-1$ define
\begin{equation}\label{eq:kappa}
\kappa_N=\frac{C_N\,q^{2-2N}}{\poch{q}{Q}{N-1}^2},\qquad
\rho_{N,r}=\frac{\poch{q^{2N-2r}}{Q}{N+r}^2\,\poch{-q^{2N-2r}}{Q}{N+r}}
{\poch{-q^{6N-2r-2}}{Q}{N+r}} .
\end{equation}

\begin{lemma}\label{lem:centre}
For $N\ge1$ and $-N\le r\le N-1$,
\begin{equation}\label{eq:centre}
h_N(N+r)=\kappa_N\,b_r\,\rho_{N,r}.
\end{equation}
Moreover,
\begin{equation}\label{eq:Sb}
S(q)=\frac{\poch{q}{Q}{\infty}^2J_2}{2J_4^4}\sum_{r\in\Z}b_r .
\end{equation}
\end{lemma}

\begin{proof}
\emph{The identity \eqref{eq:centre}.} At $r=-N$ the left side is
$h_N(0)=C_N$. On the right, $\rho_{N,-N}=1$ (all its products are empty) and
$b_{-N}=q^{2N-2}\poch{q}{Q}{N-1}^2$ by \eqref{eq:bneg}, so the right side is
also $C_N$. It remains to compare the ratios of consecutive terms. Put
$x=q^{2N}$ and $v=q^{2r}$. From \eqref{eq:hN},
\[
\frac{h_N(N+r+1)}{h_N(N+r)}=\frac{Q(1-Qv/x)^2(1+Qv/x)}{(1+Q^2v/x^3)(1-q^3v)^2},
\]
and from \eqref{eq:br} and \eqref{eq:kappa},
\[
\frac{b_{r+1}\rho_{N,r+1}}{b_r\rho_{N,r}}
=\frac{Q^2v^2}{(1-q^3v)^2}\cdot\frac{(1-x/(Qv))^2(1+x/(Qv))}{1+x^3/(Q^2v)} .
\]
After the inner fractions are cleared, both ratios equal
\[
\frac{Q\,(x-Qv)^2(x+Qv)}{(x^3+Q^2v)(1-q^3v)^2},
\]
and induction on $r$ gives \eqref{eq:centre}.

\emph{Bounds.} For each fixed $r$, the first arguments $\pm q^{2N-2r}$ and
$-q^{6N-2r-2}$ in \eqref{eq:kappa} tend to $0$ as $N\to\infty$, so
$\rho_{N,r}\to1$. Uniformly in $-N\le r\le N-1$,
\[
0\le\rho_{N,r}\le\poch{-q^2}{Q}{\infty},
\]
because $2N-2r\ge2$, the factors $1-q^{2N-2r}Q^j$ lie in $(0,1]$, the factors
$1+q^{2N-2r}Q^j$ are at most $1+q^2Q^j$, and the denominator factors are at
least $1$. Next let $B_\infty=\lim_{N\to\infty}B_N$. Letting $N\to\infty$ in
\eqref{eq:Bn} and using $\poch{Q}{Q^2}{\infty}=J_2/J_4$ and
$\poch{-Q}{Q}{\infty}=J_4/J_2$,
\[
B_\infty=\frac{\poch{q}{Q}{\infty}^4J_2}{J_4^4},\qquad
\kappa_N\longrightarrow\frac{B_\infty}{2\poch{q}{Q}{\infty}^2};
\]
for the second limit note that $C_N\sim xB_N/(2q^2)$ as $x=q^{2N}\to0$. In
particular $(\kappa_N)_{N\ge1}$ is bounded. Finally, $b_r\le
q^{2r(r+1)}/\poch{q}{Q}{\infty}^2$ for $r\ge0$, and $b_{-j-1}\le q^{2j}$ for
$j\ge0$ by \eqref{eq:bneg}; so $\sum_r b_r$ converges absolutely.

\emph{The limit.} Extend $r\mapsto h_N(N+r)$ by zero outside
$-N\le r\le N-1$. The bounds just proved give
$0\le h_N(N+r)\le(\sup_M\kappa_M)\,\poch{-q^2}{Q}{\infty}\,b_r$, a summable
majorant that does not depend on $N$. By \cref{lem:watson} and dominated
convergence,
\[
S(q)=\lim_{N\to\infty}\sum_{r}h_N(N+r)=\frac{B_\infty}{2\poch{q}{Q}{\infty}^2}\sum_{r\in\Z}b_r,
\]
which is \eqref{eq:Sb}.
\end{proof}

\begin{proof}[Proof of \cref{thm:main}]
\emph{Real $q$.} Let $0<q<1$. By \eqref{eq:bneg}, the first sum in
\cref{fact:LO} is $\sum_{r\le-2}b_r=\sum_{r<0}b_r-1$, and after the shift
$j=r+1$ the second sum is $1+\sum_{r\ge0}b_r$. Hence \cref{fact:LO} says
\begin{equation}\label{eq:bU}
\sum_{r\in\Z}b_r=-\frac{U(q)}{qJ_2}.
\end{equation}
Next split $U(q)$ by the parity of $r$. The odd terms $r=2s+1$ give
$-q^{4s^2+4s+1}/(1-q^{4s+1})$. In the even terms $r=2s$ replace $s$ by $-s$;
since $1/(1-q^{-4s-1})=-q^{4s+1}/(1-q^{4s+1})$, they give the same expression.
Therefore
\begin{equation}\label{eq:UL}
U(q)=-2q\sum_{s\in\Z}\frac{q^{4s(s+1)}}{1-q^{4s+1}}=-2qL(q).
\end{equation}
All rearrangements are allowed because the series converge absolutely.
Substituting \eqref{eq:bU} and \eqref{eq:UL} in \eqref{eq:Sb} gives
$S(q)=\poch{q}{Q}{\infty}^2L(q)/J_4^4$, which is \eqref{eq:main}.

\emph{Complex $q$.} Fix $0<t<1$. On the disc $|q|\le t$, each finite product in
$B_n$ is bounded in absolute value by a constant independent of $n$ (for
example $|\poch{q}{Q}{n}|\le\poch{-t}{t^2}{\infty}$), and each denominator
factor is bounded away from zero (for example
$|\poch{Q^2}{Q^2}{2n}|\ge\poch{t^4}{t^4}{\infty}$). The same holds for the extra
factors $1+q^{8n}$, $1+q^{2n+1}$, $1+q^{4n+2}$ in $a_n$, and $|P(q,q^{2n})|$ is
bounded. Hence $a_n=O_t(t^{2n})$, and $S$ converges absolutely and locally
uniformly on the unit disc. For the right side, rewrite the terms with $r<0$
of $L$ as in \eqref{eq:UL}:
\[
L(q)=\sum_{r\ge0}\frac{q^{4r(r+1)}}{1-q^{4r+1}}-\sum_{r\ge1}\frac{q^{4r^2-1}}{1-q^{4r-1}} .
\]
This form converges locally uniformly on the unit disc and shows that $L$ is
holomorphic there, including at $q=0$, with $L(0)=1$; the products
$\poch{q}{q^2}{\infty}$ and $J_4$ are holomorphic and nonzero there. Both sides of
\eqref{eq:main} are therefore holomorphic on the unit disc and agree on
$(0,1)$, so they agree everywhere by the identity theorem. Finally $S(0)=a_0(0)=1$
by \eqref{eq:a0}.
\end{proof}

\subsection*{Relation with Au's Example~XXI}
Let $\SXXI(q)$ be the series on the left of \cite[(3.13)]{Au}, whose summand is
\[
\frac{q^{2n}\poch{q}{Q}{n}^6}{(1+q^{2n+1})^3\poch{Q}{Q}{2n}^3}
\big(3+3q^{2n+1}-q^{4n}+2q^{4n+2}-3q^{6n+1}-3q^{8n+2}-q^{10n+3}\big).
\]
Au evaluated $\SXXI$ in \cite[(3.13)]{Au}; rewritten with \eqref{eq:UL} in the
present notation, that evaluation reads $\SXXI(q)=2J_1^4L(q)/J_2^8$. (We have
also confirmed this identity directly, as a power series modulo $q^{220}$.) Since $\poch{q}{Q}{\infty}=J_1/J_2$, comparing with
\cref{thm:main} gives the following.

\begin{corollary}\label{cor:XXI}
For $0<|q|<1$,
\[
S(q)=\frac{J_2^6}{2J_1^2J_4^4}\,\SXXI(q).
\]
\end{corollary}

Thus the Watson series and Example~XXI share the same mock modular part
$L(q)$; they differ by an infinite product, that is, by a modular form.
